\section{Quadratic jet cost for diagonal cyclic products}
\label{sec:quadratic-cost}

We prove a bound for diagonal products of the canonical cyclic cover.  The
constant in the bound depends only on an upper bound for normalized section
orders on the original fourfold.  In particular, it is independent of the
order of the cover.

The diagonal rank-deficit and determinant-order strategy adapts the
Frobenius construction of \cite[Section~9]{OpenAIAbundance}.
The cyclic many-factor weights, estimates, and constants needed here
are supplied below; the abundance conclusion of that reference is not
an input to this argument.

\begin{proposition}\label{prop:quadratic-cost}
Let $X$ be a normal projective canonical fourfold over $\C$ with
$K_X\sim_{\Q}0$, and let $r$ be the least positive integer such that
$rK_X\sim0$.  Let $\pi:Y\to X$ be its canonical cyclic cover, with group
$\mu_r$.  Fix $C_1\in\Q_{>0}$, and choose $b_0\in\Q_{>0}$ satisfying
\begin{equation}\label{eq:frob-choice-b0}
 3(2b_0)^4<\frac14,
 \qquad b_0(C_1+1)<\frac34.
\end{equation}
Suppose that $L$ is an ample rational Cartier divisor on $X$ such that
\[
 1\le L^4\le2,
 \qquad \gamma(L;X)\le C_1,
\]
where $\gamma$ is as in Definition~\ref{def:orders}.  For each integer $t\ge2$,
set
\[
 Z_t=Y^t/\mu_{r,\mathrm{diag}},
 \qquad
 \theta_t:Z_t\longrightarrow X^t,
 \qquad
 P_t=\theta_t^*\left(\sum_{i=1}^t\operatorname{pr}_i^*L\right).
\]
Then $P_t$ is ample and, at a very general point $z$ of $Z_t$,
\begin{equation}\label{eq:frob-cost}
 \varepsilon(P_t;z)\le\frac{5t^2}{b_0}.
\end{equation}
\end{proposition}

Suppose the asserted curve bound fails. The resulting surplus of
ordinary jets will produce, after reduction to characteristic $p$, a
map of vector bundles that has full rank at a general product point.
On the diagonal its rank is at most one quarter of the target rank.
A nonzero maximal minor must consequently vanish to high order there.
The section-order bound downstairs will give a strictly smaller upper
bound for that same minor at one diagonal point.

The next three lemmas supply the estimates in this comparison. The
truncated-power slope bound controls the graded pieces on the
Frobenius diagonal. The fixed flag converts those slopes into a bound
for their spaces of sections, hence for the diagonal rank. The
external-product order lemma bounds the minor by the sum of its slot
orders and rules out a larger order caused by cancellation.
We first establish these estimates on vector bundles.  For a vector bundle $E$ on a smooth
projective curve, write $\mu_{\min}(E)$ and $\mu_{\max}(E)$ for the smallest
and largest slopes in its Harder--Narasimhan filtration.

\begin{lemma}\label{lem:frob-truncated-slopes}
Let $C$ be a smooth projective integral curve of genus $g$ over an
algebraically closed field $k$ of characteristic $p>0$.  Let $V$ be a vector
bundle of rank $n$ with $\mu_{\min}(V)\ge0$, and put
\[
 G=\max\{0,2g-2\},\qquad u=n(p-1),\qquad c_2=n(n-1)G.
\]
Let $A_\bullet(V)$ denote the symmetric algebra of $V$ modulo the ideal
locally generated by the $p$-th powers of its linear generators.  Thus, in
a local frame, $A_\bullet(V)$ is the graded algebra
$\OO_C[v_1,\ldots,v_n]/(v_1^p,\ldots,v_n^p)$.  Its graded pieces are locally
free, it is zero in degrees above $u$, and
\begin{equation}\label{eq:frob-truncated-upper-slope}
 \mu_{\max}\bigl(A_j(V)\bigr)
 \le (p-1)\deg(V)+c_2
 \qquad (0\le j\le u).
\end{equation}
\end{lemma}

\begin{proof}
We use the canonical-connection approach to Frobenius instability;
see Langer \cite[\S2, Corollaries~2.4 and~6.2]{Langer2004}.
We give the estimate and tensor argument explicitly to retain the
constants needed for the many-factor diagonal.
Let $F_C$ denote absolute Frobenius.  For every vector bundle $U$ of rank
$n$, we first claim that
\begin{equation}\label{eq:frob-one-step-slope}
 \mu_{\min}(F_C^*U)
 \ge p\mu_{\min}(U)-(n-1)G.
\end{equation}
The pullback $F_C^*U$ carries the connection given by componentwise
differentiation in pulled-back frames: the transition matrices have
$p$-th-power entries.  Write the strictly decreasing slopes of its
Harder--Narasimhan factors as $\nu_1>\cdots>\nu_a$.  If some successive gap
is greater than $G$, take the last such gap, and let $S$ be the saturated
filtration step immediately before that gap.  Then
\[
 \mu_{\min}(S)>
 \mu_{\max}\bigl((F_C^*U)/S\bigr)+\deg(\Omega_C^1).
\]
Consequently the second fundamental map
\[
 S\longrightarrow ((F_C^*U)/S)\otimes\Omega_C^1
\]
vanishes, so $S$ is preserved by the connection.

Here this invariance gives Frobenius descent directly.  At the function
field $K=k(C)$, choose a matrix whose columns form
a basis of the subspace $S_K$ and whose rows in a suitable index set form
the identity matrix.  Differentiating its columns gives vectors in
$S_K\otimes\Omega^1_{K/k}$.  Their identity rows are zero, so these
differentiated columns are all zero.  Every matrix entry therefore lies
in $K^p$.  Indeed $[K:K^p]=p$ for a one-variable function field over a
perfect field, and the nonzero derivation $d:K\to\Omega^1_{K/k}$ has
kernel $K^p$.  Extracting roots of the matrix entries produces a subspace
of $U_K$.  Let $S_0\subset U$ be its saturation.  Since $C$ is smooth,
$U/S_0$ is locally free.  Frobenius is flat, so $F_C^*S_0$ is saturated
in $F_C^*U$.  It has the same generic fibre as $S$, hence equals $S$.

Put $Q=(F_C^*U)/S$.  By descent it is the Frobenius pullback of a locally
free quotient of $U$, and therefore
$\mu(Q)\ge p\mu_{\min}(U)$.  By the choice of the last large gap, all
successive Harder--Narasimhan slope gaps of $Q$ are at most $G$.  Its
average slope exceeds its smallest slope by at most $(n-1)G$.
Moreover $\mu_{\min}(Q)=\mu_{\min}(F_C^*U)$.  These observations prove
\eqref{eq:frob-one-step-slope}.  If there is no gap greater than $G$, the
same argument applies with $S=0$ and $Q=F_C^*U$.

Iteration, using $\mu_{\min}(V)\ge0$, gives for every $e\ge1$
\begin{equation}\label{eq:frob-iterated-slope}
 \mu_{\min}\bigl((F_C^e)^*V\bigr)
 \ge -(n-1)G\frac{p^e-1}{p-1}.
\end{equation}
We next deduce a tensor estimate without assuming tensor semistability in
characteristic $p$.  For each $e$, choose a line bundle $T_e$ of integral
degree $d_e$ such that
\[
 2g-\mu_{\min}\bigl((F_C^e)^*V\bigr)
 \le d_e<
 2g+1-\mu_{\min}\bigl((F_C^e)^*V\bigr).
\]
The bundle $V_e=(F_C^e)^*V\otimes T_e$ is globally generated.  In fact,
after subtracting any point its minimum slope is at least $2g-1$, so
Serre duality and the slope criterion for a nonzero homomorphism give
$H^1(C,V_e(-x))=0$; evaluation at $x$ is therefore surjective.

Let $Q_i$ be any positive-rank torsion-free quotient of $V^{\otimes i}$.
Its pullback, twisted by $T_e^{\otimes i}$, is a quotient of
$V_e^{\otimes i}$, hence is globally generated and has nonnegative
degree.  Thus
\[
 p^e\mu(Q_i)+i d_e\ge0.
\]
Divide by $p^e$ and let $e$ tend to infinity.  The upper bound for $d_e$
and \eqref{eq:frob-iterated-slope} imply
\begin{equation}\label{eq:frob-tensor-minimum}
 \mu(Q_i)\ge-\frac{i(n-1)G}{p-1}.
\end{equation}
Since $A_i(V)$ is a quotient of $V^{\otimes i}$, the same lower bound
holds for its minimum slope.  For $0\le i\le u$, it is at least $-c_2$.

Multiplication in the truncated algebra is a perfect pairing
\[
 A_j(V)\otimes A_{u-j}(V)\longrightarrow A_u(V).
\]
In a frame, each monomial pairs with its unique complementary exponent
monomial, with coefficient one.  The top line has transition character
$(\det V)^{p-1}$: this follows for diagonal changes of frame and for
elementary changes of frame, which generate the general linear group.
It follows that
\[
 A_j(V)^*\otimes(\det V)^{p-1}\simeq A_{u-j}(V).
\]
Its minimum slope is at least $-c_2$.  Slope duality now gives
\eqref{eq:frob-truncated-upper-slope}.
\end{proof}

\begin{lemma}\label{lem:frob-flag-sections}
Let $W$ be a smooth projective fourfold over an algebraically closed
field, and let $H$ be an ample rational Cartier divisor.  Fix a smooth
complete intersection flag, cut successively by divisors from
$|l_1H|$, $|l_2H|$, and $|l_3H|$, ending in a smooth integral curve $C$;
here each $l_iH$ is an integral ample divisor.  Write
$d_i=l_iH\cdot C$.  If $E$ is a vector bundle of rank $e$ and
$\mu_{\max}(E|_C)\le M$ for a real number $M\ge0$, then
\begin{equation}\label{eq:frob-flag-sections}
 h^0(W,E)\le e(M+1)\prod_{i=1}^3\left(1+\frac{M}{d_i}\right).
\end{equation}
\end{lemma}

\begin{proof}
Successive applications of the divisor exact sequences along the flag
bound $h^0(W,E)$ by
\begin{equation}\label{eq:frob-flag-sum}
 \sum_{k_1,k_2,k_3\ge0}
 h^0\left(C,E|_C\otimes
 \OO_C\left(-\sum_{i=1}^3 k_i l_iH\right)\right).
\end{equation}
To justify the iterations, at each flag stage and with the preceding
twists fixed, sufficiently negative ample twists have no sections.  A
vector bundle on that stage embeds into a finite direct sum of copies
of a sufficiently positive ample line bundle, by dualizing a global
generation map for its dual.  This gives the required vanishing and
allows each iteration of the exact sequence to terminate.  Enlarging
the resulting finite sums to \eqref{eq:frob-flag-sum} can only increase
the bound.

A summand vanishes when $\sum_i k_i d_i>M$, because its maximum slope is
then negative.  Each summand is at most $e(M+1)$: evaluation at
$\lfloor M\rfloor+1$ distinct points is injective, since subtracting
these points makes the maximum slope negative even before the additional
nonpositive twists.  There are at most
$\prod_i(1+M/d_i)$ possible triples with $\sum_i k_i d_i\le M$.
This proves \eqref{eq:frob-flag-sections}.
\end{proof}

\begin{lemma}\label{lem:frob-external-orders}
For $1\le i\le t$, let $V_i$ be an integral projective variety over an
algebraically closed field, let $x_i\in V_i$ be a smooth point, and let
$M_i$ be a line bundle with $H^0(V_i,M_i)\ne0$.  Put
\[
 a_i=\max_{0\ne s\in H^0(V_i,M_i)}\ord_{x_i}(s).
\]
Every nonzero section $s$ of the external product
$\bigotimes_i\operatorname{pr}_i^*M_i$ satisfies
\[
 \ord_{(x_1,\ldots,x_t)}(s)\le\sum_{i=1}^t a_i.
\]
\end{lemma}

\begin{proof}
Each space of sections is finite dimensional, and its order filtration
is separated; the numbers $a_i$ are therefore finite.  Choose a basis
adapted to this filtration in each slot.  The initial forms of basis
elements having a fixed order are linearly independent in that degree
of the associated graded local ring.  By the K\"unneth formula, every
section of the external product is a linear combination of the tensor
products of these bases.  In the least total degree occurring in a
nonzero combination, group its initial forms by their slot multidegrees.
Different multidegrees cannot cancel.  Within each multidegree the
products are independent, since they are tensor products of independent
initial forms in disjoint sets of smooth parameters.  The section's
order is consequently the least total degree of one of its nonzero
basis coefficients, and this is at most $\sum_i a_i$.
\end{proof}

\begin{proof}[Proof of Proposition~\ref{prop:quadratic-cost}]
The map $\theta_t$ is finite, so $P_t$ is ample.  Fix $X,L,r$, and $t$;
all auxiliary data below may depend on these choices.  Set
\begin{equation}\label{eq:frob-bq}
 b=\frac{b_0}{t},\qquad q=2tr.
\end{equation}
Suppose, for a contradiction, that
$\varepsilon(P_t;z)>5t/b$ at a smooth very general point $z$, whose image
in every $X$-slot lies in $X_{\mathrm{reg}}$.

\emph{Jets and the canonical torsor.}
Choose a rational number $c$ with
\[
 5t<c<b\varepsilon(P_t;z).
\]
On the blow-up $\sigma:\widetilde Z_t\to Z_t$ at $z$, with exceptional
divisor $E_z$, the rational divisor $\sigma^*(bP_t)-cE_z$ is ample.  To
see this, choose a larger coefficient below
$b\varepsilon(P_t;z)$ for which the corresponding divisor is nef, and
interpolate it with the ample divisor obtained at a sufficiently small
positive coefficient.  Serre vanishing, applied to sufficiently large
divisible multiples of $\sigma^*(bP_t)-cE_z$, makes the restriction map
to the corresponding thick exceptional divisor surjective.  Pullback
identifies sections of $k_0bP_t$ with sections of its pullback, by
normality of $Z_t$.  At the smooth point $z$, ordinary jets of order
strictly less than $k_0c$ inject into the sections on $k_0cE_z$.
It follows, after choosing $k_0$ large and divisible, that
\begin{equation}\label{eq:frob-initial-jets}
 \begin{gathered}
 Q=k_0bL\text{ is integral Cartier},\\
 H^0(Z_t,k_0bP_t)\longrightarrow
 \OO_{Z_t,z}/\mathfrak m_z^{5tk_0+1}
 \text{ is surjective},
 \end{gathered}
\end{equation}
where a local frame is understood in the displayed evaluation map.

Fix a section $\eta$ trivializing $rK_X$.  Over $X_{\mathrm{reg}}$,
the cover $Y\to X$ is the torsor of frames $\tau$ of the canonical line
with $\tau^r=\eta$.  Over tuples in this open locus the diagonal action
is free, and $Y^t\to Z_t$ is etale.  Choose a lift $\widetilde z$ of $z$.
Pulling back \eqref{eq:frob-initial-jets} preserves jet surjectivity at
$\widetilde z$.  The pulled-back sections are diagonal invariants.
By the K\"unneth formula and decomposition into characters, they are
linear combinations of pure tensors whose slot factors have the form
\begin{equation}\label{eq:frob-eigen-tensors}
 \begin{gathered}
 \tau^{-m_i}s_i,
 \qquad
 s_i\in H^0\bigl(X_{\mathrm{reg}},\OO(Q+m_iK_X)\bigr),
 \\
 0\le m_i<r,
 \qquad
 \sum_i m_i\equiv0\pmod r.
 \end{gathered}
\end{equation}
Indeed multiplication of an eigenvector by the corresponding power of
the tautological frame makes it invariant and hence descending to the
indicated line bundle.

Take a smooth projective resolution $u:W\to X$, isomorphic over
$X_{\mathrm{reg}}$, and write
\[
 K_W=u^*K_X+A,\qquad A\ge0\text{ exceptional}.
\]
We use the same letters $L,Q$ for their pullbacks to $W$.
Every $s_i$ in \eqref{eq:frob-eigen-tensors} extends to a section of
$Q+m_iK_W$.  In fact, its effective zero divisor extends by closure on
the normal variety $X$; this divisor is rational Cartier, since its
class is $Q+m_iK_X$.  Pulling it back and adding $m_iA$ gives the
effective divisor of the corresponding rational section on $W$.
Similarly $\eta$ extends to
\[
 \eta_W\in H^0(W,rK_W),\qquad \Div(\eta_W)=rA.
\]
Let $U=u^{-1}(X_{\mathrm{reg}})$.  On the torsor of roots of $\eta_W$
over $U$, a fixed finite list of the tensors
\eqref{eq:frob-eigen-tensors}, now using these sections on $W$, spans
the jets in \eqref{eq:frob-initial-jets} at the chosen lift.

\emph{A polarization, a flag, and a movable test.}
Fix an ample divisor $H_{\mathrm{amp}}$ on $W$.  Since $L$ is nef and
$K_WL^3=0$, a sufficiently small positive rational $\delta$ makes
$H=L+\delta H_{\mathrm{amp}}$ ample and ensures
\begin{equation}\label{eq:frob-polarization}
 H^4\le3,
 \qquad LH^3\le H^4,
 \qquad (q+t-1)K_WH^3\le b_0H^4.
\end{equation}
The equality $K_WL^3=0$ follows from $K_X\sim_{\Q}0$ and the
exceptionality of $A$.  The canonical class of $W$ is pseudoeffective.
By \cite[Theorem~2.1]{CampanaPaun2015}, $\Omega_W^1$ is generically
semipositive for this polarization.  Apply
\cite[Theorem~6.1]{MehtaRamanathan1982} to its finitely many
Harder--Narasimhan factors.  We can choose a smooth complete
intersection flag, cut successively by sufficiently positive integral
multiples $l_iH$, $1\le i\le3$, ending in a smooth integral curve $C$
such that
\begin{equation}\label{eq:frob-flag-nonnegative}
 \mu_{\min}(\Omega_W^1|_C)\ge0.
\end{equation}
The general curve avoids the codimension-two loci where the factors
or their filtration quotients fail to be locally free.  Their
restricted slopes are the original slopes multiplied by the same
positive factor $l_1l_2l_3$, so the restriction theorem gives
\eqref{eq:frob-flag-nonnegative}.

Choose also a very general point $x\in U$ at which the all-degree bound
defining $\gamma(L;X)\le C_1$ holds.  Let
$\pi_x:\widehat W\to W$ be its blow-up, with exceptional divisor $J$,
and put
\begin{equation}\label{eq:frob-test-divisor}
 D^+=bL+(q+1)K_W,\qquad d_0=b(C_1+1).
\end{equation}
The divisor $D^+$ is big.  In sufficiently divisible degrees its
section spaces agree with those of $bL$, after multiplication by a
fixed section whose zero divisor is exceptional over $X$.  This is
because $rK_X\sim0$, $A$ is effective and exceptional, and
$u_*\OO_W(E)=\OO_X$ for any effective integral exceptional divisor
$E$ on this resolution.  At $x$ these fixed sections have order zero.
Thus every nonzero section of a sufficiently divisible multiple
$mD^+$ has order at most $mbC_1$ at $x$.

It follows that $\pi_x^*D^+-d_0J$ is not pseudoeffective.  Otherwise,
interpolating it with the big divisor $\pi_x^*D^+$ would make
$\pi_x^*D^+-dJ$ big for a rational $d$ with $bC_1<d<d_0$.  An effective
sufficiently divisible multiple would give a section of a multiple of
$D^+$ with an excluded order at $x$.
The projective algebraic duality theorem of
\cite[Theorem~2.2, arXiv version]{BDPP2013} therefore supplies a strongly movable
class $\beta$ on $\widehat W$ such that
\begin{equation}\label{eq:frob-movable-test}
 D^+\cdot\beta<d_0J\cdot\beta,
 \qquad J\cdot\beta>0.
\end{equation}
Here and subsequently divisors on $W$ are pulled back when paired with
$\beta$.  More precisely, we may take
$\beta=v_*(A_1A_2A_3)$ for a smooth projective birational model
$v:T\to\widehat W$ and very ample divisors $A_i$ on $T$.  The dual cone
is the closure of the cone generated by such classes; strict negativity
on that cone yields negativity on one generator.  Also
$D^+\cdot\beta\ge0$, since $L$ is nef and $K_W$ is pseudoeffective.
The second inequality in \eqref{eq:frob-movable-test} follows from the
first and $d_0>0$.  In particular,
\begin{equation}\label{eq:frob-test-nonnegative}
 L\cdot\beta\ge0,\qquad K_W\cdot\beta\ge0.
\end{equation}

\emph{Reduction of the fixed data.}
Spread the preceding finite algebraic data over an integral finitely
generated $\Z$-subalgebra of $\C$.  This includes the smooth varieties,
the flag, the points and the blow-up, the morphism $v$ and its very ample
divisors, suitable integral multiples of the rational line bundles,
the sections used in \eqref{eq:frob-eigen-tensors}, the section
$\eta_W$, and the torsor and its marked lift over $U$.  After shrinking
the base, the relevant fibres of $W,C,T$, and $\widehat W$ are smooth,
projective, and geometrically integral; the testing morphism remains
dominant; and all required ampleness and intersection data persist.
The fixed finite jet evaluation remains surjective: it is a map of
finite locally free modules obtained by restriction to a fixed-order
infinitesimal neighbourhood of a section of a smooth morphism, and
surjectivity is open.  Invert $r$ and the finitely many denominators
used in the data.

We can also retain \eqref{eq:frob-flag-nonnegative} on every geometric
fibre of this smaller base.  To see this without imposing any
Frobenius semistability condition, choose a fixed relative ample twist
that makes the family $\Omega_W^1|_C$ relatively globally generated.
For every positive-rank quotient, its degree is then bounded below
uniformly, in terms of its rank and this twist.  Hence there are only
finitely many possible negative degrees and ranks, and therefore
finitely many Hilbert polynomials, for quotients of negative degree.
The corresponding relative Quot schemes are projective.  None has a
point on the geometric generic fibre: such a point would persist after
extension to the original complex field and contradict
\eqref{eq:frob-flag-nonnegative}.  This also excludes quotients with
torsion and negative degree, since their torsion-free quotients have
still smaller degree.  Remove the images of these finitely many Quot
schemes.  The remaining geometric fibres have
$\mu_{\min}(\Omega_W^1|_C)\ge0$.

There are geometric fibres in arbitrarily large characteristics $p$:
the base is nonempty, of finite type, and dominant over $\Spec\Z$.
We take such fibres over algebraic closures of residue fields and keep
the preceding notation.  The genus of $C$ and all the intersection
numbers used above are constant.  The class $\beta$ continues to pair
nonnegatively with every effective divisor: pull the divisor back
under the dominant morphism $v$ and intersect with $A_1A_2A_3$.
This argument applies also after a base-field twist.  Thus only the
fixed test, and not pseudoeffective-cone duality, is being specialized.

\emph{Full rank away from the diagonal.}
In one of these reductions, write $F:W\to W'$ for relative Frobenius,
where a prime denotes base-field twist.  It is finite flat of degree
$p^4$.  Put
\begin{equation}\label{eq:frob-ranks}
 \begin{gathered}
 N=\lfloor p/k_0\rfloor,
 \qquad B=NQ,
 \qquad \mathcal E=F_*\OO_W(B),\\
 \qquad s=p^4=\rk(\mathcal E),
 \qquad R=s^t.
 \end{gathered}
\end{equation}
Products of $N$ of the fixed jet tensors span jets through degree
$5tk_0N$ at the chosen lift.  Indeed every monomial of degree at most
$5tk_0N$ is a product of $N$ monomials of degree at most $5tk_0$.
Choose representatives with these leading terms from the original
jet-spanning space.  Their products give a triangular spanning set for
the jet filtration, in every characteristic.  For all sufficiently
large $p$,
\[
 5tk_0N\ge4t(p-1).
\]
The local monomial inclusion used here is the one in
\cite[proof of Proposition~2.12, equation~(2.8)]{MustataSchwede2014},
with smooth dimension $4t$. The products therefore span on the thick point defined by the $p$-th
powers of $4t$ smooth parameters.  Etaleness identifies the completed
local ring at the torsor lift with the completed local ring at its
tuple in $U^t$.  Under this identification the thick point is the
fibre of $F^t:W^t\to(W')^t$ over the Frobenius image of that tuple.

In a product of $N$ pure tensors, let $m_i'$ denote the new exponent
in slot $i$.  Then
\[
 0\le m_i'\le N(r-1)<pr,
 \qquad \sum_i m_i'\equiv0\pmod r.
\]
Since $p$ is prime to $r$, choose integers $a_i$ with
\[
 r\le a_i\le2r-1,\qquad pa_i\equiv m_i'\pmod r.
\]
Their sum is a multiple of $r$ and is at most $2tr-t<q$.  Increase
$a_1$ by the nonnegative multiple $q-\sum_i a_i$ of $r$.  We now have
\begin{equation}\label{eq:frob-padded-weights}
 a_i\ge0,\qquad \sum_i a_i=q,\qquad pa_i\ge m_i'.
\end{equation}
Multiply the corresponding slot section of $B+m_i'K_W$ by
$\eta_W^{(pa_i-m_i')/r}$.  This gives a section of $B+pa_iK_W$, and
division by $\tau^{pa_i}$ gives exactly its previous torsor expression.
In a local canonical frame, the coefficient of $\tau$ is a unit.  Its
$p$-th power is a nonzero constant on the thick Frobenius fibre, since
the $p$-th powers of all local parameters vanish there.

It follows that the projection-formula map
\begin{equation}\label{eq:frob-column-map}
 \bigoplus_{\substack{a_i\ge0\\\sum_i a_i=q}}
 \left(\bigotimes_{i=1}^tH^0(W,\OO_W(B+pa_iK_W))\right)
 \otimes
 \OO_{(W')^t}\left(-\sum_{i=1}^t a_i\operatorname{pr}_i^*K_{W'}\right)
 \longrightarrow\mathcal E^{\boxtimes t}
\end{equation}
has rank $R$ at one point, hence at the generic point.  The products
just constructed give the spanning values, up to the nonzero constant
frame factors.  By finite pushforward and base change these values are
the fibre of the target.  The identity used in the projection formula
is $F^*\OO_{W'}(K_{W'})\simeq\OO_W(pK_W)$; it is an identity of line
bundles, not the pullback map on differential forms.  In compatible
local frames its transitions are $p$-th powers.

\emph{The loss of rank on the diagonal.}
On the diagonal $W'\subset(W')^t$, all the source line bundles in
\eqref{eq:frob-column-map} equal $\OO_{W'}(-qK_{W'})$.  Let
\[
 \Delta_F=\underbrace{W\times_{W'}\cdots\times_{W'}W}_{t\ \mathrm{factors}},
 \qquad h:\Delta_F\to W',\qquad g:\Delta_F\to W
\]
be the common structure map and the first projection, respectively.
Both are finite.  The line bundle
\begin{equation}\label{eq:frob-diagonal-line}
 \mathcal L=\OO_{\Delta_F}\left(
       \sum_{i=1}^t\operatorname{pr}_i^*B
       +pq\operatorname{pr}_1^*K_W\right)
\end{equation}
satisfies
\[
 h_*\mathcal L\simeq
 (\mathcal E^{\boxtimes t})|_{W'}\otimes\OO_{W'}(qK_{W'}).
\]
This is finite base change from $F^t$, followed by the projection
formula; all the line bundles $\operatorname{pr}_i^*\OO_W(pK_W)$
are the pullback of $\OO_{W'}(K_{W'})$.  The columns of
\eqref{eq:frob-column-map}, after twisting, thus restrict to global
sections of $h_*\mathcal L$.  At every point of the diagonal their rank
is at most $h^0(\Delta_F,\mathcal L)$.

The two-factor diagonal filtration and its truncated symmetric pieces
are the canonical filtration of
\cite[Theorem~3.7]{Sun2008} and
\cite[Definition~3.1, Remark~3.2, and Corollary~3.5]{KitadaiSumihiro2008}.
We use the same local construction in $t-1$ difference slots, and give
the resulting ranks and determinant comparison explicitly.
Filter $\mathcal L$ by powers of the ideal of the reduced diagonal
$W\subset\Delta_F$ and push forward by $g$.  The associated pieces are
\begin{equation}\label{eq:frob-diagonal-graded}
 \begin{gathered}
 \mathcal G_j=\OO_W(tB+pqK_W)\otimes A_j(V),
 \qquad V=(\Omega_W^1)^{\oplus(t-1)},\\
 0\le j\le u_0=4(t-1)(p-1).
 \end{gathered}
\end{equation}
For completeness, in etale smooth coordinates the other slots add
$4(t-1)$ difference variables with their $p$-th powers zero and no
other relations.  This follows from the Cartesian relative Frobenius
square for an etale chart.  Modulo squares the differences transform
by ordinary differentials, giving $V$; the associated graded algebra
is therefore intrinsically $A_\bullet(V)$.  The line bundle contributes
its restriction $\OO_W(tB+pqK_W)$ to each piece.  In particular, with
$e_j=\rk A_j(V)$,
\begin{equation}\label{eq:frob-diagonal-section-sum}
 h^0(\Delta_F,\mathcal L)\le\sum_{j=0}^{u_0}h^0(W,\mathcal G_j),
 \qquad \sum_{j=0}^{u_0}e_j=p^{4(t-1)}=s^{t-1}.
\end{equation}

Set $n=4(t-1)$, $G=\max\{0,2g(C)-2\}$, and $c_2=n(n-1)G$.
By the retained inequality \eqref{eq:frob-flag-nonnegative},
$\mu_{\min}(V|_C)\ge0$.  Lemma~\ref{lem:frob-truncated-slopes} gives,
uniformly in $j$ and $p$,
\begin{align}
 \mu_{\max}(\mathcal G_j|_C)
 &\le
 \bigl(tB+pqK_W+(p-1)(t-1)K_W\bigr)\cdot C+c_2\notag\\
 &\le M_p:=2b_0p\,l_1l_2l_3H^4+c_2.
 \label{eq:frob-Mp}
\end{align}
Indeed $BH^3\le pbLH^3$, $K_WH^3\ge0$, and
\eqref{eq:frob-polarization} bounds the two contributions by
$pb_0H^4$ each, before multiplication by $l_1l_2l_3$.

Put $a=l_1l_2l_3H^4$ and $d_i=l_iH\cdot C=l_i a$.  Applying
Lemma~\ref{lem:frob-flag-sections} to each $\mathcal G_j$ yields
\begin{align}
 h^0(W,\mathcal G_j)
 &\le e_j(M_p+1)\prod_{i=1}^3(1+M_p/d_i)\notag\\
 &=e_jp^4\bigl((2b_0)^4H^4+O(p^{-1})\bigr).
 \label{eq:frob-uniform-sections}
\end{align}
The error is uniform in $j$, since $c_2,a,d_i$ were fixed before $p$
varied.  Explicitly, the leading coefficient is
$(2b_0)^4a^4/\prod_i d_i=(2b_0)^4H^4$.
Summing \eqref{eq:frob-uniform-sections}, using
\eqref{eq:frob-diagonal-section-sum}, $H^4\le3$, and the strict first
inequality in \eqref{eq:frob-choice-b0}, gives, for all sufficiently
large $p$,
\begin{equation}\label{eq:frob-diagonal-rank}
 h^0(\Delta_F,\mathcal L)\le R/4.
\end{equation}
Thus the rank of \eqref{eq:frob-column-map} is at most $R/4$ at every
point of the diagonal.

\emph{The determinant contradiction.}
Choose $R$ columns of \eqref{eq:frob-column-map} with nonzero
determinant, and write $a_i^{(\nu)}$ for the weight of slot $i$ in
column $\nu$.  Their determinant is a nonzero section of the external
product of the line bundles
\[
 \mathcal Q_i=(\det\mathcal E)^{\otimes R/s}
     \otimes\OO_{W'}\left(\sum_{\nu=1}^R a_i^{(\nu)}K_{W'}\right).
\]
Let $\lambda_i=R^{-1}\sum_{\nu=1}^R a_i^{(\nu)}$; then
$0\le\lambda_i\le q$.
Identifying numerical classes through base-field twist, we have
\begin{equation}\label{eq:frob-determinant-class}
 \frac{c_1(\mathcal Q_i)}{R}
 =\frac{B}{p}
   +\left(\frac{p-1}{2p}+\lambda_i\right)K_W.
\end{equation}
The first-Chern coefficient is the usual Frobenius direct-image
coefficient; compare \cite[Lemma~4.2]{Sun2008}. Here only its
numerical divisor class is needed. To check it directly, finite base change identifies
$F^*\mathcal E$ with the pushforward from $W\times_{W'}W$ of the line
bundle pulled back from its second slot.  The same truncated-diagonal
filtration used in \eqref{eq:frob-diagonal-graded}, now with two slots,
has pieces $\OO_W(B)\otimes A_j(\Omega_W^1)$.
Across all $p^4=s$ truncated monomials, the total exponent of each of
the four variables is $s(p-1)/2$.  Taking determinants gives
\[
 c_1(F^*\mathcal E)=sB+\frac{s(p-1)}2K_W.
\]
Relative Frobenius pullback multiplies divisor classes by $p$ after
the base-twist identification.  Since the determinant of
$\mathcal E^{\boxtimes t}$ contributes
$(\det\mathcal E)^{R/s}$ in each slot, this proves
\eqref{eq:frob-determinant-class}.

Use now the base twist of the blow-up of $x$ and of the test class
$\beta$.  By $B/p\le bL$ on this test,
\eqref{eq:frob-test-nonnegative}, and $\lambda_i\le q$, the normalized
class in \eqref{eq:frob-determinant-class} pairs with the test to at
most $D^+\cdot\beta$.  Let $s_i$ be any nonzero section of
$\mathcal Q_i$, and let $m_i=\ord_{x'}(s_i)$.  Its divisor has effective
strict transform of class
$\pi_{x'}^*c_1(\mathcal Q_i)-m_iJ'$ on the twisted blow-up.  The
nonnegativity of the test on effective divisors and
\eqref{eq:frob-movable-test} imply
\[
 m_i(J\cdot\beta)
 \le c_1(\mathcal Q_i)\cdot\beta
 \le R D^+\cdot\beta
 <R d_0(J\cdot\beta).
\]
Here the pairings have been identified with their untwisted numerical
values.  As $J\cdot\beta>0$, every such section has order strictly less
than $Rd_0$ at $x'$.

The determinant section is nonzero, so the K\"unneth formula ensures
that each slot section space is nonzero.  By
Lemma~\ref{lem:frob-external-orders}, its order at
$(x',\ldots,x')$ is therefore strictly less than
\begin{equation}\label{eq:frob-order-upper}
 Rt d_0=R b_0(C_1+1)<3R/4.
\end{equation}
On the other hand, \eqref{eq:frob-diagonal-rank} bounds the rank of its
column matrix at this point by $R/4$.  Trivialize the source and target
bundles locally, and perform constant invertible row operations so
that at least $R-\lfloor R/4\rfloor$ rows vanish at the point.  Every
entry of these rows lies in the maximal ideal.  Each term of the
determinant is consequently in its
$(R-\lfloor R/4\rfloor)$-th power.  The determinant has order at least
$3R/4$, contradicting \eqref{eq:frob-order-upper}.

This excludes the assumed Seshadri inequality and proves
\eqref{eq:frob-cost}.  All varieties, line bundles, flags, points, and
testing data were fixed before the characteristic grew.  Their
dependence on $X,r,t$ therefore affects only how large the
characteristic must be, whereas the constant $5/b_0$ in the conclusion
depends only on $C_1$.
\end{proof}
